\chapter{Teorijske osnove prevodilaca}
\label{ch:teorija}

\napomena{Teorijsko poglavlje, 10--12 strana. Uputstvo: „ovdje će trebati istraživanja”.
Glavni izvori: \textcite{aho2006} (pogl.\ 1--6), \textcite{beck1997} (pogl.\ 1--2),
\textcite{cooper2023}, \textcite{nystrom2021}.}

\section{Vrste prevodilaca}
\label{sec:vrste-prevodilaca}
\napomena{Kompajler, interpreter, asembler, (dis)asembler, kros-prevodilac; jednoprolazni
i višeprolazni prevodioci.}

\section{Struktura prevodioca}
\label{sec:struktura-prevodioca}
\napomena{Faze: front-end (analiza) i back-end (sinteza); tabela simbola kao zajednička
struktura. Dijagram faza.}

\section{Leksička analiza}
\label{sec:teorija-leksicka}
\napomena{Tokeni, lekseme, obrasci; regularni izrazi i konačni automati;
ručno pisani skener naspram generatora (lex/flex).}

\section{Sintaksna analiza}
\label{sec:teorija-sintaksna}
\napomena{Kontekstno-slobodne gramatike, BNF/EBNF, stabla parsiranja;
silazno (rekurzivno spuštanje) i uzlazno parsiranje.}

\section{Semantička analiza}
\label{sec:teorija-semanticka}
\napomena{Provjere značenja: definisanost simbola, opsezi vrijednosti, tipovi operanada;
kod asemblera: nedefinisane/duplirane labele, neposredne vrijednosti van opsega.}

\section{Generisanje mašinskog koda}
\label{sec:teorija-generisanje}
\napomena{Kod asemblera: kodiranje instrukcija u bitska polja; dvoprolazni postupak
i reference unaprijed \parencite{beck1997}; jednoprolazni asembler sa „backpatching”-om.}

\section{Obrada grešaka}
\label{sec:teorija-greske}
\napomena{Vrste grešaka (leksičke, sintaksne, semantičke), lokacija greške,
strategije oporavka (panic mode, nastavak na sljedećoj liniji).}
